\section*{Computing the Laplacian of the Spectral Density Estimator}

We estimate the marginal density using a truncated spectral expansion of the Laplace--Beltrami operator:
\[
\hat f(x)
=
\frac{1}{n}\sum_{i=1}^n K(x,X_i,\ell),
\qquad
K(x,y,\ell)
=
\sum_{j\le \ell} \phi_j(x)\phi_j(y),
\]
where $\{\phi_j\}$ are eigenfunctions of the Laplace--Beltrami operator $\Delta$:
\[
\Delta \phi_j = \lambda_j \phi_j.
\]

\subsection*{Key Observation}

Since the eigenfunctions diagonalize $\Delta$, we have
\[
\Delta_x K(x,y,\ell)
=
\sum_{j\le \ell} \lambda_j \phi_j(x)\phi_j(y).
\]

Therefore, the Laplacian of the density estimator is
\[
\boxed{
\Delta \hat f(x)
=
\frac{1}{n}
\sum_{i=1}^n
\sum_{j\le \ell}
\lambda_j \phi_j(x)\phi_j(X_i).
}
\]

Thus, computing $\Delta \hat f$ does \emph{not} require taking second derivatives explicitly.  
It suffices to multiply each spectral coefficient by its eigenvalue.

\subsection*{Example: Circle $S^1$}

On $S^1$, the eigenfunctions are
\[
\phi_k(x)=\frac{1}{\sqrt{2\pi}}e^{ikx},
\qquad
\Delta \phi_k = -k^2 \phi_k.
\]

Hence,
\[
\Delta \hat f(\phi)
=
\sum_{k=-M}^{M} (-k^2)\,\hat c_k e^{ik\phi}.
\]

\subsection*{Example: Sphere $S^2$}

For zonal expansions using Legendre polynomials $P_m$,
\[
\Delta P_m(\langle x,y\rangle)
=
- m(m+1) P_m(\langle x,y\rangle).
\]

Thus each degree-$m$ term is multiplied by $-m(m+1)$ when computing the Laplacian.

\bigskip

\noindent
In summary, when using a Laplace--Beltrami spectral estimator, the Laplacian is obtained by simply multiplying each retained eigencomponent by its eigenvalue.